% answers_short.tex —— 2011 年·上海卷·理科 · 简短答案速查
% 由 tools/make_short_answers.py 自 answers.tex 的 \gksolhead/\gkanswer 提取，
% 并与 2011/2011_数学真题_上海理_答案要点.md 逐题比对后生成，勿手改（改源 MD 后重跑）。

\begin{gknotice}
\textbf{使用说明}：本节只给最终答案便于快速核对，完整推导见解析版。答案为编写者独立推导，\textbf{非官方标准答案}。
\end{gknotice}

\gkpart{填空题}
\begin{itemize}[leftmargin=2.6em]
  \item[\textbf{1.}] $f^{-1}(x)=\dfrac1x+2$（$x\neq0$）。
  \item[\textbf{2.}] $(0,1)$。
  \item[\textbf{3.}] $16$。
  \item[\textbf{4.}] $(-\infty,0)\cup\left[\dfrac12,+\infty\right)$。
  \item[\textbf{5.}] $\arctan\dfrac12$（即 $\arccos\dfrac{2\sqrt5}5$）。
  \item[\textbf{6.}] $\sqrt6$。
  \item[\textbf{7.}] $\dfrac{\sqrt3\pi}3$。
  \item[\textbf{8.}] $\dfrac{2+\sqrt3}4$。
  \item[\textbf{9.}] $2$。
  \item[\textbf{10.}] $6$。
  \item[\textbf{11.}] $\dfrac{15}2$。
  \item[\textbf{12.}] $0.985$。
  \item[\textbf{13.}] $[-16,12]$。
  \item[\textbf{14.}] $\sqrt3$。
\end{itemize}

\gkpart{选择题}
\begin{center}\begin{tabular}{ccccc}
\toprule
题号 & 15 & 16 & 17 & 18 \\
\midrule
答案 & D & A & B & D \\
\bottomrule
\end{tabular}\end{center}

\gkpart{解答题}
\begin{itemize}[leftmargin=2.6em]
  \item[\textbf{19.}] $z_2=4+2\mathrm i$。
  \item[\textbf{20.}] (1) $a,b>0$ 时在 $\mathbf R$ 上单调递增；$a,b<0$ 时在 $\mathbf R$ 上单调递减；(2) $a>0>b$ 时 $x<\log_{2/3}\dfrac{-2b}{a}$；$a<0<b$ 时 $x>\log_{3/2}\dfrac{-a}{2b}$
  \item[\textbf{21.}] (1) 证明见解析；(2) $2$
  \item[\textbf{22.}] (1) $c_1=9$，$c_2=11$，$c_3=12$，$c_4=13$；(2) 证明见解析；(3) $c_{4k+1}=6k+9$，$c_{4k+2}=6k+11$，$c_{4k+3}=6k+12$，$c_{4k+4}=6k+13$（$k=0,1,2,\cdots$）
  \item[\textbf{23.}] (1) $\sqrt5$；(2) $4+\pi$；(3) ① $\Omega$ 为 $y$ 轴；② $x=0\ (y\geqslant0)$，$x=\dfrac{y^2}4\ (-2\leqslant y<0)$，$x=-y-1\ (y<-2)$；③ $x\leqslant0,\ y\leqslant0$；$y=x\ (0<x\leqslant1)$；$y=\dfrac{x^2+1}2\ (1<x\leqslant2)$；$x=\dfrac{2y+3}4\ \left(y>\dfrac52\right)$
\end{itemize}
